% 第6章 循环控制结构（教材第6章 6.1--6.8）
\lstset{language=C}

\section{循环控制结构与循环语句}

实际应用中的许多问题都涉及\textbf{重复执行}一些操作，如级数求和、穷举或迭代求解。若重复处理的次数是已知的，称为\textbf{计数控制的循环}（Counter Controlled Loop）；若次数未知、由给定条件控制，称为\textbf{条件控制的循环}（Condition Controlled Loop）。两者都需要用循环结构实现。

循环结构有两种类型：

\begin{itemize}
  \item \textbf{当型循环}：先判断条件 P，成立（真）时反复执行 A 操作，直到条件不成立（假）时结束——\textbf{先测后做}；
  \item \textbf{直到型循环}：先执行 A 操作，再判断条件 P，成立则继续反复执行 A，直到条件不成立——\textbf{先做后测}。
\end{itemize}

这个差别直接决定了循环体\textbf{最少执行几次}：当型可能一次都不执行，直到型至少执行一次。

C 语言提供 \texttt{while}、\texttt{do-while}、\texttt{for} 三种循环语句来实现循环结构。被重复执行的语句序列称为\textbf{循环体}（Body of Loop）。

\subsection{while 语句}

\texttt{while} 属于\textbf{当型}循环：

\begin{lstlisting}
while (循环控制表达式)
{
    语句序列;      /* 循环体 */
}
\end{lstlisting}

执行过程：

\begin{enumerate}
  \item 计算循环控制表达式的值；
  \item 若为真，执行循环体语句，然后返回步骤 1；
  \item 若为假，退出循环，执行循环体后面的语句。
\end{enumerate}

循环控制表达式在\textbf{执行循环体之前}测试，因此若第一次判断就为假，循环体一次也不执行。

\subsection{do-while 语句}

\texttt{do-while} 属于\textbf{直到型}循环，注意结尾的分号不能少：

\begin{lstlisting}
do
{
    语句序列;      /* 循环体 */
} while (循环控制表达式);
\end{lstlisting}

与 \texttt{while} 不同的是，\texttt{do-while} 的循环控制表达式在\textbf{执行循环体之后}测试，因此\textbf{循环体内的语句至少被执行一次}。这正是“先输入再校验”这类场景需要的语义：至少让用户输一次，才能判断输入是否合法。

\subsection{for 语句}

\texttt{for} 也属于当型循环，是 C 程序中使用频率最高的循环语句：

\begin{lstlisting}
for (初始化表达式; 循环控制表达式; 增值表达式)
{
    语句序列;      /* 循环体 */
}
\end{lstlisting}

三个表达式的分工：

\begin{itemize}
  \item \textbf{初始化表达式}：为循环控制变量赋初值，决定循环的起始条件；
  \item \textbf{循环控制表达式}：控制循环继续执行的条件，值为真（非 0）时继续，否则结束循环；
  \item \textbf{增值表达式}：每执行完一次循环体后对循环控制变量做一次变化，决定循环何时结束。
\end{itemize}

每轮（包括第一轮）循环体执行之前都要测试一次循环控制条件，每轮循环体执行完后都要执行一次增值表达式。\texttt{for} 语句与下面形式的 \texttt{while} 语句等价：

\begin{lstlisting}
初始化表达式;
while (循环控制表达式)
{
    语句序列;
    增值表达式;
}
\end{lstlisting}

\textbf{三条书写规则}：

\begin{enumerate}
  \item 三个表达式之间的分隔符是\textbf{分号}，有且仅有两个分号，不能多也不能少；
  \item 循环控制表达式一般不要省略，若省略则表示循环条件恒真（死循环）。若已在 \texttt{for} 前面赋过初值，初始化表达式可省略；若已在循环体中改变了循环控制变量，增值表达式可省略。因此 \texttt{for (; cond; )}、\texttt{for (i=0; cond; )} 都是合法的；
  \item 即使循环体只有一条语句，也建议用花括号括起来。否则日后往循环体里加语句时，只有第一条在循环内，会造成逻辑错误。
\end{enumerate}

\begin{lstlisting}
/* 三种语句实现同一件事：累加 1 到 n */
sum = 0;
for (i = 1; i <= n; i++)  sum += i;

sum = 0;  i = 1;
while (i <= n) { sum += i; i++; }

sum = 0;  i = 1;
do { sum += i; i++; } while (i <= n);
\end{lstlisting}

\section{计数控制的循环}

计数控制循环的典型应用是\textbf{累加求和}与\textbf{累乘求积}。以计算 $1+2+\cdots+n$ 为例，如果为每个加数定义一个变量，$n$ 未知时就无法确定变量个数，$n$ 很大时又浪费空间。改用循环后，只需三个变量：循环控制变量 \texttt{i}、累加和 \texttt{sum}、循环次数 \texttt{n}。

\textbf{（1）累加变量必须初始化为 0，累乘变量必须初始化为 1。}这是最容易被忽略、后果最严重的一步。删掉 \texttt{sum = 0} 后，\texttt{sum} 里是残留的垃圾值，累加结果会变成 \texttt{4199309} 之类的乱数——因为“在垃圾值的基础上加”。

\textbf{（2）\texttt{sum = sum + i} 为什么能累加？}因为赋值号右侧的 \texttt{sum} 执行的是\textbf{读}操作，左侧的 \texttt{sum} 执行的是\textbf{写}操作：先把 \texttt{sum} 的旧值读出来，与 \texttt{i} 相加，再写回 \texttt{sum}。所以每轮循环都“在上一次的基础上继续累积”。

\textbf{（3）循环变量的出口值要注意。}对 \texttt{for (i = 1; i <= n; i++)}，循环结束时 \texttt{i} 的值是 $n+1$ 而不是 $n$，因为最后一轮循环体执行完后还要执行一次增值表达式才做判断。

\textbf{（4）找通项规律是计数循环的核心。}当被累加的各项不是简单的 $i$，而是 $i^{2}$、$\frac{1}{i}$、$(-1)^{i}\frac{1}{2i-1}$ 等形式时，关键是写出\textbf{第 $i$ 项}的表达式，并尽量用\textbf{递推}的办法由上一项算出下一项，而不是每轮都从头计算：

\begin{lstlisting}
/* 累乘求阶乘：term 表示当前项，每轮由上一项递推 */
double term = 1.0, result = 0.0;
for (i = 1; i <= n; i++) {
    term = term * i;       /* 第 i 项 = 第 i-1 项 * i */
    result = result + term;
}
\end{lstlisting}

\textbf{（5）循环次数的确定。}判断循环执行多少次，看的是“初值、终值、步长”三件事：\texttt{for (i = 1; i <= n; i++)} 执行 $n$ 次；\texttt{for (i = 0; i < n; i++)} 也执行 $n$ 次；而 \texttt{for (i = 1; i < n; i++)} 只执行 $n-1$ 次。边界上的等号是“差一错误”（off-by-one）的根源。

\section{嵌套循环}

一个循环体内又包含另一个完整的循环结构，称为\textbf{嵌套循环}。外层循环每执行一次，内层循环就要完整地执行一遍，因此总执行次数是\textbf{各层次数之积}。以打印图形为例：

\begin{lstlisting}
for (i = 1; i <= rows; i++) {        /* 外层：控制行数 */
    for (j = 1; j <= i; j++)         /* 内层：控制本行字符数 */
        putchar('*');
    putchar('\n');                    /* 每行结束换行 */
}
\end{lstlisting}

嵌套循环要注意三点：

\begin{enumerate}
  \item \textbf{内外层循环变量不能同名}。同名会互相覆盖，导致内层破坏外层的计数，常常直接造成死循环或次数错乱；
  \item \textbf{内层循环的初值必须每轮重新设置}。若把 \texttt{j} 的初始化写到外层循环之外，第二轮起 \texttt{j} 已经是旧值，内层可能一次都不执行；
  \item \textbf{花括号要成对}，尤其在内层循环体只有一条语句时，漏写花括号会让“每行换行”的语句跑进内层，输出变成一列。
\end{enumerate}

\section{条件控制的循环}

当循环次数事先未知、由某个条件决定何时结束时，用条件控制的循环。典型场景是“反复读入数据，直到满足某条件为止”：

\begin{lstlisting}
/* 猜数游戏：反复猜，直到猜中 */
do {
    printf("Please guess: ");
    scanf("%d", &guess);
    count++;
    if (guess > target)      printf("too big\n");
    else if (guess < target) printf("too small\n");
} while (guess != target);   /* 猜中即结束 */
printf("You got it in %d times\n", count);
\end{lstlisting}

这类循环有两条设计要点：

\begin{itemize}
  \item \textbf{循环控制条件必须在循环体内被改变}，否则一旦进入就是死循环。上面的 \texttt{guess} 每轮由 \texttt{scanf} 更新，\texttt{count} 每轮递增，条件终会成立；
  \item 若“至少要做一次”才谈得上判断（如先读入才能判断是否结束），用 \texttt{do-while}；若“可能一次都不做”，用 \texttt{while}。
\end{itemize}

\section{流程的转移控制}

\texttt{goto}、\texttt{break}、\texttt{continue} 和 \texttt{return} 是 C 语言中用于控制流程转移的跳转语句，其中 \texttt{return} 用于从函数返回值，在第 7 章介绍。

\subsection{goto 语句}

\texttt{goto} 是\textbf{无条件转向}语句，既可向下跳转，也可往回跳转：

\begin{lstlisting}
goto 语句标号;
...
语句标号:  语句;
\end{lstlisting}

语句标号（Label）代表跳转的目标位置，必须用合法标识符表示，命名规则与变量名相同。\texttt{goto} 通常与 \texttt{if} 联合使用，构成“满足条件才跳”的形式。

良好的编程风格要求\textbf{少用、慎用 \texttt{goto}}，尤其不要用往回跳转的 \texttt{goto}，也不要用它制造永远不会被执行的死代码。结构化程序设计规定：程序中\textbf{尽量不使用多于一个 \texttt{goto} 语句标号}，且只允许在“单入口单出口”的模块内\textbf{向前}跳转，不允许回跳。

\texttt{goto} 真正合适的情形只有两种：一是从多层嵌套循环中一次性跳出，二是跳向共同的出口位置做统一的错误处理。

\subsection{break 语句}

\texttt{break} 除用于退出 \texttt{switch} 结构外，还可用于 \texttt{while}、\texttt{do-while}、\texttt{for} 的循环体中。执行到 \texttt{break} 时，\textbf{循环立即终止}，从循环语句后的第一条语句继续执行。

它是“有条件跳转”，且跳转位置被唯一限定为\textbf{紧接循环语句后的第一条语句}。正因为位置固定，它不需要标号，也就避免了因标号过多导致的流程混乱。

\textbf{关键限制}：在嵌套循环中，\texttt{break} \textbf{只对包含它的最内层循环起作用}，不能一次跳出多重循环。要跳出 $k$ 层，只能逐层 \texttt{break}，或者用 \texttt{goto}。

\subsection{continue 语句}

\texttt{continue} 与 \texttt{break} 都用于循环内部控制，但效果不同：遇到 \texttt{continue} 时，程序\textbf{跳过本次循环中后面尚未执行的语句，开始下一次循环}——只结束本次循环，并不终止整个循环。

\begin{lstlisting}
for (i = 1; i <= 5; i++) {
    scanf("%d", &n);
    if (n < 0) {
        printf("negative, skipped\n");
        continue;            /* 不显示负数，但继续读下一个 */
    }
    printf("n = %d\n", n);
}
\end{lstlisting}

\textbf{一个必须记住的差异}：\texttt{for} 循环里的 \texttt{continue} 会\textbf{执行增值表达式}再判断条件；而\texttt{while} 循环里的 \texttt{continue} 会\textbf{直接跳去判断条件}，若增值语句写在 \texttt{continue} 之后就被跳过，从而造成死循环。所以“把 \texttt{for} 改写成 \texttt{while}”在有 \texttt{continue} 时二者\textbf{并不等价}——这是教科书反复提醒的一点。

\subsection{穷举法编程实例}

\textbf{穷举}（Exhaustion）就是通过尝试问题的所有可能来得到最终答案。以“韩信点兵”为例：按 1 至 5 报数末位为 1，按 1 至 6 报数末位为 5，按 1 至 7 报数末位为 4，按 1 至 11 报数末位为 10，求至少有多少兵。设兵数为 $x$，则需

\[
x\%5=1 \ \&\&\ x\%6=5 \ \&\&\ x\%7=4 \ \&\&\ x\%11=10
\]

从 1 开始逐个试验，第一个使条件成立的 $x$ 即为所求：

\begin{lstlisting}
int x = 1;
while (!(x % 5 == 1 && x % 6 == 5 && x % 7 == 4 && x % 11 == 10))
    x++;
printf("x = %d\n", x);
\end{lstlisting}

穷举法有三个必须交代清楚的问题：

\begin{enumerate}
  \item \textbf{枚举范围要覆盖解}。若把上界写成 \texttt{x < 3000}，而真正的解在范围之外，程序将“一无所获”。偶然得到正确答案不算对，那是瞎猫碰到死耗子；
  \item \textbf{要有终止机制}。若去掉上界限制，虽然一定能找到解，但会得到不止一个解，而且找完第一个解后循环还会继续，最终陷入死循环。要么用 \texttt{break} 在找到后立即退出，要么在条件里体现“至少”的含义；
  \item \textbf{枚举量要尽量小}。可以借助数学性质剪枝（如把步长从 1 改为 11 的倍数），在解空间很大时这一点决定程序能否跑完。
\end{enumerate}

\section{本章扩充内容}

\subsection{结构化程序设计的核心}

结构化程序设计的核心是：\textbf{任何复杂问题都只用顺序、选择、循环三种基本结构来编程实现}，并且每个模块保持“单入口、单出口”。由此派生出几条具体规定：

\begin{itemize}
  \item 限制和避免使用 \texttt{goto}，最多一个标号，只允许向前跳；
  \item 采用自顶向下、逐步求精的模块化方法分解问题（第 7 章）；
  \item 程序结构的好坏不由“有没有 \texttt{goto}”决定，避免 \texttt{goto} 只是手段而非目的。
\end{itemize}

\subsection{常用的程序调试与排错方法}

调试的目标不是“用工具乱翻”，而是\textbf{先想清楚，再验证}。几条实用做法：

\begin{itemize}
  \item \textbf{插入打印语句}往往比单步跟踪更高效：在关键位置打印中间变量，观察值是否符合预期。排错语句留在程序中，可反复观察；单步执行则转瞬即逝；
  \item \textbf{增量测试}：扩充程序时不断测试新加入的代码及它与原有可运作单元的联系，把待调试的代码限制在很小的范围内，从而缩小查错区域；
  \item \textbf{改错三步}：先想清楚修改这一行会不会引发其他缺陷；乘胜追击，改正所有同类 bug；改完立即做\textbf{回归测试}（对以前修过的 bug 重新测试），确认没有引入新问题；
  \item \textbf{重视错误现象}。错误结果常常提供线索：问题为什么发生？以前是否遇到类似情况？哪几条语句刚被修改过？引起错误的输入有什么特殊之处？
\end{itemize}

\subsection{类型溢出}

C 语言提供的每种类型都有取值范围，向其中赋超过范围的值就会产生\textbf{类型溢出}（Type Overflow），得到错误结果。典型例子是“计算自然数的立方和，直到立方和大于等于 1000000”：

\begin{lstlisting}
short i, sum = 0;              /* 错误：short 只有 2 字节 */
for (i = 1; ; i++) {
    sum = sum + i * i * i;
    if (sum >= 1000000) { printf("count = %d\n", i); break; }
}
\end{lstlisting}

\texttt{short} 只占 2 字节，取值范围是 $-32768 \sim 32767$，而 1000000 早已超出，\texttt{sum} 会溢出成负数，条件永远不成立，于是\textbf{死循环}。正确结果应为 45 项，把变量改为 \texttt{int} 或 \texttt{long} 即可。

溢出还有两种更隐蔽的形式：

\begin{itemize}
  \item \textbf{中间结果溢出}：即便最终结果在范围内，中间乘积也可能越界。例如 \texttt{n * n} 中 \texttt{n} 为 \texttt{int} 时，$n$ 稍大就会溢出，此时应先强转为 \texttt{double} 再乘；
  \item \textbf{整数除法丢精度}：计算 $\frac{n^{2}}{(n-1)(n+1)}$ 时，若 \texttt{n} 为 \texttt{int}，\texttt{n*n / ((n-1)*(n+1))} 是整数除法，小数部分被舍去，必须写成 \texttt{(double)(n*n) / ((double)(n-1)*(n+1))}。
\end{itemize}

教训是：\textbf{预先估算运算结果的可能取值范围，选用取值范围足够大的类型}；把高位的平台移植到低位平台时尤其要重新检查。

\section{本章知识点小结}

\begin{itemize}
  \item 循环分当型（\texttt{while}、\texttt{for}，先测后做）与直到型（\texttt{do-while}，先做后测，至少执行一次）；
  \item \texttt{for} 的三个表达式由两个分号分隔；初始化、增值可省略，条件省略即死循环；
  \item 累加变量初值为 0，累乘变量初值为 1；\texttt{sum = sum + i} 的读写次序是理解累加的关键；
  \item 嵌套循环总次数为各层次数之积；内外层变量不能同名，内层初值每轮必须重置；
  \item \texttt{goto} 可任意跳转但破坏结构；\texttt{break} 终止整个循环且只跳出最内层；\texttt{continue} 只结束本次循环；
  \item \texttt{for} 中的 \texttt{continue} 会执行增值表达式，\texttt{while} 中不会，二者不一定等价；
  \item 穷举法要保证枚举范围覆盖解、有终止机制，并尽量缩小枚举空间；注意类型溢出与整数除法丢精度。
\end{itemize}

\section{本章常见错误小结}

\begin{table}[htbp]
\centering
\caption{第 6 章常见错误与纠正}
\small
\begin{tabularx}{\textwidth}{@{}lXl@{}}
\toprule
易错写法 & 错在哪里 & 纠正 \\
\midrule
累加前未 \texttt{sum = 0} & 在垃圾值上累加，输出乱数 & 循环前初始化 \\
\texttt{for (i = 1; i <= n; i++);} & 圆括号后多写分号，循环体变成空语句 & 去掉分号 \\
\texttt{while (cond);} & 同上，且容易变成死循环 & 去掉分号 \\
循环体中未改变控制变量 & 条件永远成立，死循环 & 循环体内更新变量 \\
\texttt{for (i=1, j=1; i<n; i++, j++)} & 三个表达式之间用逗号而非分号 & 用两个分号分隔三部分 \\
\texttt{do \{...\} while (cond)} & 漏写结尾分号，编译错误 & \texttt{while (cond);} \\
\texttt{break} 想跳出两层循环 & \texttt{break} 只作用于最内层 & 逐层 \texttt{break} 或用 \texttt{goto}/\texttt{return} \\
\texttt{short sum} 累加大数 & 类型溢出，结果为负，循环不终止 & 用 \texttt{int}/\texttt{long} \\
\bottomrule
\end{tabularx}
\end{table}
